\documentclass[10pt,a4paper]{amsart}
\usepackage[margin=19mm]{geometry}
\usepackage{amsmath,amssymb,mathtools}
\usepackage[scr=rsfs,cal=euler]{mathalfa}
\usepackage{xfrac}
\usepackage[unicode,hidelinks,bookmarks=false]{hyperref}
\newcommand{\ie}{i.e.\ }
\newcommand{\cf}{cf.\ }
\newcommand{\resp}{respectively\ }
\newcommand{\Subsection}{\subsection}
\providecommand{\nobreakdash}{\nobreak\hbox{-}}
\setlength{\emergencystretch}{2em}
\multlinegap0pt
\usepackage{amssymb}
\usepackage{enumerate}
\usepackage{tikz}
\usepackage{float}
\usepackage[shortlabels]{enumitem}
\newtheorem{theorem}{Theorem}
\title{Distribution of signless Laplacian eigenvalues and degree sequence}
\author{Saieed Akbari, M. Darougheh, L. S. de Lima, D. Tracina}
\date{Source: arXiv:2605.27405v1 (CC BY 4.0)}
\begin{document}
\maketitle

\noindent\emph{Source and licence.} Distribution of signless Laplacian eigenvalues and degree sequence. Version arXiv:2605.27405v1 (CC BY 4.0), distributed by arXiv under the Creative Commons Attribution 4.0 International licence, http://creativecommons.org/licenses/by/4.0/, as recorded in the arXiv licence metadata for this exact version and shown on the abstract page https://arxiv.org/abs/2605.27405v1. Full licence notice: \texttt{licenses/arxiv-2605.27405-CC-BY-4.0.txt}.\par\smallskip

\noindent\textit{Editorial scope.} Counted unit: the article's theorem theorem (source0.tex lines 1806--1808). The article's complete original proof of it is delivered with it (source0.tex lines 1810--1875, one \texttt{proof} environment of 66 lines). No mathematics is written, repaired or re-derived: the statement and the proof are the article's own lines, in the article's own notation, and no retained line is substituted or re-flowed.  No further source-local context is needed: every cross-reference of every retained line resolves inside the two delivered spans.  Nothing was omitted from any retained span, so the package carries no omission record.  No retained line contains a citation, so the wrapper carries no bibliography and no bibliography entry.  No retained line contains a figure environment, an image-inclusion command or drawing code, so no image is delivered and the package declares no asset dependency.  Editorial formatting only, in addition: an AMS \texttt{amsart} preamble that restates the article's own declarations and macros used by the retained lines, namely the environments \texttt{theorem} on separate counters, together with the standard packages those lines need.  The retained environments print in this document as Theorem 1 in order of appearance, whereas the article prints the same items with the numbering of its own full document.  The complete native source of the article as distributed in the arXiv e-print for this exact version is bundled unchanged as \texttt{sources/201489/source0.tex}.
\begin{theorem}
There is no graph of order $n \ge 3$ such that ${m_G}[0,{d_3}] = 1$.
\end{theorem}

\begin{proof}
By computer computation we checked the assertion for all graphs up to $5$ vertices.
Thus assume that $n\ge 6$. For each vertex $v_i\in V(G)$ let $d_G(v_i)=d_i$ for $1\le i\le n$.
By the pigeonhole principle there exist at least two vertices of degree at most $d_3$ such that either both are neighbors of $v_3$ or both are non-neighbors of $v_3$. Without loss of generality assume these vertices are $v_4$ and $v_5$.

\noindent Define vectors $x=(x_1,\dots,x_n)^T$ and $y=(y_1,\dots,y_n)^T\in\mathbb{R}^n$ by
\[
	\begin{aligned}
		x_i=
		\begin{cases}
			{1 \over \sqrt{2}} & i = 4,\\
			{-1 \over \sqrt{2}} & i = 5,\\
			0 & \text{otherwise},
		\end{cases}
		\qquad
		y_i=
		\begin{cases}
			1 & i = 3,\\
			0 & \text{otherwise}.
		\end{cases}
	\end{aligned}
	\]
Let $S=\operatorname{span}\{x,y\}$ and let $Z\in S$ with $\|Z\|=1$.
By the Courant--Fischer theorem,
\[
q_{n-1}\le \max_{\substack{Z\in S\\ \|Z\|=1}} Z^T Q Z.
\]
Write $Z=a x + b y$ with $a,b\in\mathbb{R}$; since $x\perp y$ and $\|Z\|=1$ we have $a^2+b^2=1$.
Then
\[
Z^T Q Z = a^2\, x^T Q x + b^2\, y^T Q y + 2ab\, x^T Q y.
\]

\noindent Clearly, we get
\[
y^T Q y = d_3.
\]
Notice that ${x^T}Qy = \sum\limits_{{v_i}{v_j} \in E(G)} {{x_i}{q_{ij}}{y_j}} $. Since $v_3$ is either adjacent to both $v_4$ and $v_5$, or to neither, so we obtain that ${x^T}Qy = 0.$

\noindent In addition,
\[
	\begin{aligned}
		x^T Q x &=\sum_{v_i v_j\in E(G)} (x_i+x_j)^2=
		\begin{cases}
			{{{{d_4} + {d_5}} \over 2} - 1} & v_4v_5\in E(G),\\
			{{{{d_4} + {d_5}} \over 2}} & v_4v_5\notin E(G).\\
		\end{cases}
	\end{aligned}
\]

\noindent Therefore
\[
Z^T Q Z = a^2 \alpha + b^2 d_3,
\]
where
\[
\alpha =
\begin{cases}
{{{{d_4} + {d_5}} \over 2} - 1} & v_4v_5\in E(G),\\
			{{{{d_4} + {d_5}} \over 2}} & v_4v_5\notin E(G).\\
		
\end{cases}
\]
Since $a^2+b^2=1$ and $\alpha  \le {d_3}$, the maximum of $a^2 \alpha + b^2 d_3$ occurs when $a=0$ and $b=1$.
Thus ${q_{n - 1}} \le {d_3}\ $, and the proof is complete.
\end{proof}

\end{document}
